\subsection{Neighboring operators and commuting bond products}

\begin{definition}[Neighboring operators]%
    \label{def:mpdo_sector_eta}
    \lean{MPOTensor.sectorEta}
    \leanok
    \uses{def:mpdo_sector_tensors}
    Fix outer virtual indices $\alpha_1,\beta_3$ with the chosen Hayashi
    decomposition, and let $l_k,r_k$ be the sector tensors of
    Definition~\ref{def:mpdo_sector_tensors}. For each pair of sectors
    $(k,h)$, the neighboring operator contracts the right sector tensor of
    one site with the left sector tensor of the following site over the
    virtual bond they share:
    \begin{align}
        \eta_{k,h}
         & =\sum_{\gamma}(r_k)_{\gamma}\otimes(l_h)_{\gamma}.
        \notag
    \end{align}
    It acts on the neighboring bond space $B_k^{R}\otimes B_h^{L}$. This is
    the operator displayed in
    \cite[Appendix~C.2, lines~1441--1445]{Cirac2016MPDO_arXiv}.
\end{definition}

\begin{theorem}[Neighboring bond contraction]
    \label{thm:mpdo_neighboring_bond_contraction}
    \lean{MPOTensor.physicalSlice_neighboring_contraction}
    \leanok
    \uses{def:mpdo_sector_eta}
    In the setting of Theorem~\ref{thm:mpdo_sector_factorization}, for all
    sectors $k,h$ and virtual indices $\beta_1,\alpha_3$, contracting the
    shared virtual bond of two neighboring factorized slices leaves the
    neighboring operator between the outer sector tensors:
    \begin{align}
         & \sum_{\gamma}
        \left[U_B\kappa_{\beta_1,\gamma}U_B^\dagger\right]_{(k,l_1,r_1),(k,l_1',r_1')}
        \left[U_B\kappa_{\gamma,\alpha_3}U_B^\dagger\right]_{(h,l_2,r_2),(h,l_2',r_2')}
        \notag           \\
         & \quad =
        \left[(l_k)_{\beta_1}\right]_{l_1,l_1'}
        \left[\eta_{k,h}\right]_{(r_1,l_2),(r_1',l_2')}
        \left[(r_h)_{\alpha_3}\right]_{r_2,r_2'}.
        \notag
    \end{align}
    This is the single-bond step of the identity assembling
    $\bigotimes_{n}\eta_{k_n,k_{n+1}}$ from the factorized chain in
    \cite[Appendix~C.2, lines~1446--1449]{Cirac2016MPDO_arXiv}.
\end{theorem}

\begin{proof}
    \leanok
    \uses{thm:mpdo_sector_factorization}
    Substituting the sector factorization
    (Theorem~\ref{thm:mpdo_sector_factorization}) into each slice, the
    outer sector tensors decouple from the bond sum, and the inner sum over
    the shared bond $\gamma$ reads
    \begin{align}
        \sum_{\gamma}
        [(r_k)_{\gamma}]_{r_1,r_1'}
        [(l_h)_{\gamma}]_{l_2,l_2'}
         & =
        [\eta_{k,h}]_{(r_1,l_2),(r_1',l_2')},
        \notag
    \end{align}
    by definition of the neighboring operator.
\end{proof}

\begin{definition}[Source-minimal physical-sector factorization]
    \label{def:mpdo_physical_sector_factorization}
    \lean{MPOTensor.PhysicalSectorFactorization,
        MPOTensor.PhysicalSectorFactorization.transformedPhysicalSlice,
        MPOTensor.PhysicalSectorFactorization.transformedPhysicalSlice_eq,
        MPOTensor.PhysicalSectorFactorization.sectorFinEquiv,
        MPOTensor.PhysicalSectorFactorization.sectorCoordinateTensor,
        MPOTensor.PhysicalSectorFactorization.neighborIndex_nonempty}
    \lean{MPOTensor.PhysicalSectorFactorization.sectorCoordinateTensor_apply,
        MPOTensor.PhysicalSectorFactorization.sectorCoordinateTensor_apply_same,
        MPOTensor.PhysicalSectorFactorization.sectorCoordinateTensor_apply_ne}
    \lean{MPOTensor.PhysicalSectorFactorization.transformedPhysicalSlice_apply_same,
        MPOTensor.PhysicalSectorFactorization.transformedPhysicalSlice_apply_ne}
    \leanok
    \uses{def:mpo_tensor}
    A physical-sector factorization of an MPO tensor $\mc K$ consists of a
    finite decomposition
    \begin{align}
        \C^d
         & \simeq\bigoplus_k(B_k^L\otimes B_k^R),
        \notag
    \end{align}
    whose left and right factors are nonzero,
    an isometry $U$ on the physical space, and matrix families
    $(l_k)_\beta$ and $(r_k)_\alpha$ such that every physical slice satisfies
    \begin{align}
        U\kappa_{\beta,\alpha}U^\dagger
         & =\bigoplus_k(l_k)_\beta\otimes(r_k)_\alpha.
        \notag
    \end{align}
    No positivity, trace factorization, or quantum-Markov decomposition is
    included in this datum. This is the factorization displayed in
    \cite[Appendix~C.2, lines~1381--1388]{Cirac2016MPDO_arXiv}.
\end{definition}

\begin{definition}[Algebraic neighboring contraction]
    \label{def:mpdo_minimal_neighboring_operator}
    \lean{MPOTensor.PhysicalSectorFactorization.neighboringOperator,
        MPOTensor.PhysicalSectorFactorization.neighboringOperator_apply}
    \leanok
    \uses{def:mpdo_physical_sector_factorization}
    For sectors $k,h$, contract the common virtual index of the right and
    left sector tensors:
    \begin{align}
        \eta_{k,h}
         & =\sum_a(r_k)_a\otimes(l_h)_a.
        \notag
    \end{align}
    This is an operator on $B_k^R\otimes B_h^L$. Positivity is not part of
    this definition; it is a separate hypothesis in Proposition~C.7.
\end{definition}

\begin{definition}[Matrix-weighted neighboring contraction]
    \label{def:mpdo_matrix_weighted_neighboring_operator}
    \lean{MPOTensor.PhysicalSectorFactorization.neighboringOperatorWithMatrix,
        MPOTensor.PhysicalSectorFactorization.neighboringOperatorWithMatrix_apply,
        MPOTensor.PhysicalSectorFactorization.neighboringOperatorWithMatrix_one}
    \leanok
    \uses{def:mpdo_physical_sector_factorization, def:mpdo_minimal_neighboring_operator}
    For a virtual matrix $P$, define
    \begin{align}
        \eta_{k,h}(P)
         & =\sum_{\delta,\gamma}P_{\delta,\gamma}
        (r_k)_\delta\otimes(l_h)_\gamma.
        \notag
    \end{align}
    The ordinary neighboring contraction is the case $P=I$.
\end{definition}

\begin{definition}[Virtual-matrix transport of a physical-sector factorization]
    \label{def:mpdo_physical_sector_factorization_virtual_matrices}
    \lean{MPOTensor.PhysicalSectorFactorization.ofVirtualMatrices}
    \leanok
    \uses{def:mpdo_physical_sector_factorization}
    Let $K$ and $L$ have virtual dimensions $D$ and $E$, respectively.
    Suppose that matrices $X\in M_{E\times D}(\C)$ and
    $Y\in M_{D\times E}(\C)$ satisfy the following identity for every
    physical index $i$:
    \begin{align}
        L^i & =XK^iY.
        \notag
    \end{align}
    Every physical-sector factorization of $K$ induces one of $L$ with the
    same physical decomposition.  Its virtual tensor families are
    \begin{align}
        (l_k^L)_\beta
         & =\sum_\gamma X_{\beta,\gamma}(l_k^K)_\gamma,
         &
        (r_k^L)_\alpha
         & =\sum_\delta Y_{\delta,\alpha}(r_k^K)_\delta.
        \notag
    \end{align}
    This construction abstracts the virtual corner in
    \cite[Section~2.3, lines~195--219]{Cirac2016MPDO_arXiv} and the
    physical-sector factorization in
    \cite[Appendix~C.2, lines~1381--1439]{Cirac2016MPDO_arXiv}.
\end{definition}

\begin{theorem}[Neighboring operators under virtual-matrix transport]
    \label{thm:mpdo_physical_sector_factorization_virtual_matrices_neighbor}
    \lean{MPOTensor.PhysicalSectorFactorization.ofVirtualMatrices_neighboringOperator}
    \leanok
    \uses{def:mpdo_minimal_neighboring_operator, def:mpdo_matrix_weighted_neighboring_operator, def:mpdo_physical_sector_factorization_virtual_matrices}
    The neighboring operators of the transported factorization satisfy
    \begin{align}
        \eta^L_{k,h}
         & =\eta^K_{k,h}(YX).
        \notag
    \end{align}
    This identity combines the virtual corner in
    \cite[Section~2.3, lines~195--219]{Cirac2016MPDO_arXiv} with the
    neighboring contraction in
    \cite[Appendix~C.2, lines~1441--1450]{Cirac2016MPDO_arXiv}.
\end{theorem}

\begin{proof}
    \leanok
    Replace each left factor by
    $\sum_\gamma X_{\beta,\gamma}(l_k)_\gamma$ and each right factor by
    $\sum_\delta Y_{\delta,\alpha}(r_k)_\delta$.  Contracting the common
    virtual index gives
    \begin{align}
        \sum_\beta Y_{\delta,\beta}X_{\beta,\gamma}
         & =(YX)_{\delta,\gamma},
        \notag
    \end{align}
    which proves the neighboring-operator identity.
\end{proof}

\begin{definition}[Virtual compression of a physical-sector factorization]
    \label{def:mpdo_physical_sector_factorization_virtual_compression}
    \lean{MPOTensor.PhysicalSectorFactorization.ofVirtualCompression}
    \leanok
    \uses{def:mpdo_physical_sector_factorization_virtual_matrices}
    Let $K$ and $L$ have virtual dimensions $D$ and $E$, respectively, and
    suppose that a rectangular matrix $V\in M_{D\times E}(\C)$ satisfies the
    following identity for every physical index $i$:
    \begin{align}
        L^i & =V^\dagger K^iV.
        \notag
    \end{align}
    Every physical-sector factorization of $K$ induces one of $L$ with the
    same physical decomposition.
    This construction combines the virtual corner in
    \cite[Section~2.3, lines~195--219]{Cirac2016MPDO_arXiv} with the
    physical-sector factorization in
    \cite[Appendix~C.2, lines~1381--1439]{Cirac2016MPDO_arXiv}.
\end{definition}

\begin{theorem}[Range projection of a matrix-block inclusion]
    \label{thm:block_inclusion_projection}
    \lean{Matrix.sigmaBlockInclusion_mul_conjTranspose_eq_blockProjection}
    \leanok
    For the canonical inclusion $V_s$ of a matrix block,
    $V_sV_s^\dagger$ is the projection onto that block.
\end{theorem}

\begin{proof}\leanok
    Compute the entries of the inclusion and its conjugate transpose.
\end{proof}

\begin{theorem}[Flattened sector as an isometric virtual corner]
    \label{thm:mpdo_flattened_sector_virtual_compression}
    \lean{MPSTensor.SectorDecomposition.flatBlockInclusion,
        MPSTensor.SectorDecomposition.flatBlockInclusion_isometry,
        MPSTensor.SectorDecomposition.flatBlockInclusion_mul_conjTranspose,
        MPSTensor.SectorDecomposition.toTensor_mul_flatBlockInclusion,
        MPSTensor.SectorDecomposition.flatBlockInclusion_conjTranspose_mul_toTensor,
        MPSTensor.SectorDecomposition.flatWeight_smul_flatBasis_eq_compression}
    \leanok
    \uses{def:sector_weight_data, thm:block_inclusion_projection}
    Let $A=\bigoplus_t \mu_t A_t$ be a weighted block-diagonal tensor in
    flattened virtual coordinates.
    For every block $s$, the canonical inclusion $V_s$ of that block satisfies
    \begin{align}
        V_s^\dagger V_s    & =1, \notag           \\
        V_sV_s^\dagger     & =P_s, \notag         \\
        V_s^\dagger A^iV_s & =\mu_s A_s^i. \notag
    \end{align}
    Here $P_s$ is the corresponding block projection, transported to the
    flattened coordinates.  Moreover,
    \begin{align}
        A^iV_s          & =V_s(\mu_sA_s^i), \notag         \\
        V_s^\dagger A^i & =(\mu_sA_s^i)V_s^\dagger. \notag
    \end{align}
    This is the exact blockwise virtual-corner identity from
    \cite[Section~2.3, equation \emph{II\_Aiplusk1},
        lines~195--219]{Cirac2016MPDO_arXiv}.  Appendix~C.2, Case~I,
    lines~1374--1450 and 1571--1619 is the downstream application context;
    no positivity after insertion of $P_s$ is asserted here.
\end{theorem}

\begin{proof}
    \leanok
    Let $E_s$ be the canonical inclusion of the $s$-th summand in the
    dependent direct sum, let $Q_s$ be the corresponding unflattened block
    projection, and let $e$ flatten the direct-sum coordinates.  Thus $V_s$
    is the reindexing of $E_s$, while $P_s$ is the reindexing of $Q_s$.
    Multiplicativity of reindexing transports
    \begin{align}
        E_s^\dagger E_s & =1, \notag \\
        E_sE_s^\dagger  & =Q_s
        \notag
    \end{align}
    to the two corresponding identities for $V_s$.  If
    $B_t^i=\mu_tA_t^i$, block diagonality gives
    \begin{align}
        A^iV_s          & =V_sB_s^i, \notag  \\
        V_s^\dagger A^i & =B_s^iV_s^\dagger.
        \notag
    \end{align}
    Therefore
    \begin{align}
        V_s^\dagger A^iV_s
         & =V_s^\dagger V_sB_s^i
        =B_s^i
        =\mu_sA_s^i.
        \notag
    \end{align}
\end{proof}

\begin{theorem}[Neighboring operators under virtual compression]
    \label{thm:mpdo_physical_sector_factorization_virtual_compression_neighbor}
    \lean{MPOTensor.PhysicalSectorFactorization.ofVirtualCompression_neighboringOperator}
    \leanok
    \uses{def:mpdo_minimal_neighboring_operator, def:mpdo_matrix_weighted_neighboring_operator, def:mpdo_physical_sector_factorization_virtual_compression}
    The neighboring operators of the factorization obtained by virtual
    compression satisfy
    \begin{align}
        \eta^L_{k,h}
         & =\eta^K_{k,h}(VV^\dagger).
        \notag
    \end{align}
    No positivity of the right-hand side is asserted.
    This identity combines the virtual corner in
    \cite[Section~2.3, lines~195--219]{Cirac2016MPDO_arXiv} with the
    neighboring contraction in
    \cite[Appendix~C.2, lines~1441--1450]{Cirac2016MPDO_arXiv}.
\end{theorem}

\begin{proof}
    \leanok
    \uses{thm:mpdo_physical_sector_factorization_virtual_matrices_neighbor}
    Apply the virtual-matrix transport identity with $X=V^\dagger$ and $Y=V$.
\end{proof}

\begin{theorem}[Virtual-gauge invariance of physical-sector factorizations]
    \label{thm:mpdo_physical_sector_factorization_gauge}
    \lean{MPOTensor.PhysicalSectorFactorization.ofGaugeEquiv,
        MPOTensor.PhysicalSectorFactorization.ofGaugeEquiv_eq_ofVirtualMatrices,
        MPOTensor.PhysicalSectorFactorization.ofGaugeEquiv_neighboringOperator,
        MPOTensor.PhysicalSectorFactorization.ofGaugeEquiv_neighboringOperator_posSemidef}
    \leanok
    \uses{def:mpdo_physical_sector_factorization, def:mpdo_minimal_neighboring_operator, def:mpdo_physical_sector_factorization_virtual_matrices}
    Suppose that two MPO tensors of the same virtual dimension differ by an
    invertible virtual gauge $X$.  The induced factorization is precisely the
    virtual-matrix transport by $X$ and $X^{-1}$.  Thus every physical-sector
    factorization of the first tensor induces one of the second tensor with
    the same physical decomposition and the same neighboring operators.  In
    particular, positive semidefiniteness of all neighboring operators is
    invariant under the gauge.
\end{theorem}

\begin{proof}
    \leanok
    \uses{def:mpdo_matrix_weighted_neighboring_operator, thm:mpdo_physical_sector_factorization_virtual_matrices_neighbor}
    Specialize the virtual-matrix neighboring identity to a gauge and its
    inverse.  Their product is the identity matrix, and contraction through
    the identity matrix is the ordinary neighboring operator.
\end{proof}

\begin{definition}[One-site matrix spaces of sector sets]%
    \label{def:mpdo_one_site_sector_matrix_space}
    \lean{MPOTensor.PhysicalSectorFactorization.oneSiteSectorMatrix,
        MPOTensor.PhysicalSectorFactorization.oneSiteSectorSpace}
    \leanok
    \uses{def:mpdo_physical_sector_factorization}
    For a sector $k$ and physical matrix indices $x,y$ in its summand, let
    $A_{k;x,y}\in\MN{D}$ be the corresponding entry of the
    sector-coordinate tensor. For a set $S$ of sectors, define
    \begin{align}
        \mc A_S
         & =\spn\{A_{k;x,y}:k\in S\}.
        \notag
    \end{align}
    These are the matrix spaces associated with the physical projections in
    \cite[Appendix~C.2, lines~1463--1467]{Cirac2016MPDO_arXiv}.
\end{definition}

\begin{theorem}[Directed-cut annihilation of one-site matrix spaces]%
    \label{thm:mpdo_directed_cut_one_site_spaces}
    \lean{MPOTensor.PhysicalSectorFactorization.%
        oneSiteSectorMatrix_mul_eq_zero_of_neighboringOperator_eq_zero}
    \lean{MPOTensor.PhysicalSectorFactorization.oneSiteSectorSpace_mul_eq_zero}
    \leanok
    \uses{def:mpdo_minimal_neighboring_operator, def:mpdo_one_site_sector_matrix_space}
    Let $S$ and $C$ be sets of sectors. If $\eta_{k,h}=0$ for every
    $k\in S$ and $h\in C$, then $XY=0$ for every $X\in\mc A_S$ and
    $Y\in\mc A_C$. No conclusion in the reverse product order is asserted.
    This is the directed local-orthogonality implication used in
    \cite[Appendix~C.2, lines~1463--1470]{Cirac2016MPDO_arXiv}.
\end{theorem}

\begin{proof}
    \leanok
    Expand two generating matrices in sector coordinates. Their product has
    entries
    \begin{align}
        \sum_{\gamma}
        (l_k)_\beta (r_k)_\gamma
        (l_h)_\gamma (r_h)_\alpha.
        \notag
    \end{align}
    The middle contraction is a matrix entry of $\eta_{k,h}$ and therefore
    vanishes. Bilinearity extends the conclusion from the generators to the
    two linear spans.
\end{proof}

\begin{definition}[Vertex rephasing of the sector tensors]%
    \label{def:mpdo_physical_sector_rephasing}
    \lean{MPOTensor.PhysicalSectorFactorization.rephase}
    \leanok
    \uses{def:mpdo_physical_sector_factorization, def:mpdo_minimal_neighboring_operator}
    Let $z_k\in\C$ satisfy $|z_k|=1$ for every sector. Replacing
    \begin{align}
        (l_k)_a & \longmapsto z_k(l_k)_a,
        \notag                                \\
        (r_k)_a & \longmapsto z_k^{-1}(r_k)_a
        \notag
    \end{align}
    leaves every physical-slice factorization unchanged.
\end{definition}

\begin{theorem}[Rephasing identities for the sector data]%
    \label{thm:mpdo_physical_sector_rephasing_invariance}
    \lean{MPOTensor.PhysicalSectorFactorization.rephase_neighboringOperator}
    \lean{MPOTensor.PhysicalSectorFactorization.rephase_cyclicNeighboringProduct}
    \lean{MPOTensor.PhysicalSectorFactorization.rephase_sectorVirtualMatrix}
    \lean{MPOTensor.PhysicalSectorFactorization.%
        rephase_neighboringOperator_ne_zero_iff}
    \leanok
    \uses{def:mpdo_physical_sector_rephasing, def:mpdo_minimal_neighboring_operator}
    Vertex rephasing transforms the neighboring operators by
    $\eta_{k,h}\longmapsto z_k^{-1}z_h\,\eta_{k,h}$. It leaves every sector
    virtual matrix unchanged, preserves the nonzero neighboring support, and
    leaves every complete cyclic product of neighboring operators unchanged.
    These are the rephasing identities used in
    \cite[Appendix~C.2, lines~1434--1450]{Cirac2016MPDO_arXiv}.
\end{theorem}

\begin{definition}[Vanishing right tensors on inactive sectors]
    \label{def:mpdo_physical_sector_inactive_right_zero}
    \lean{MPOTensor.PhysicalSectorFactorization.zeroRightTensorOnInactive}
    \lean{MPOTensor.PhysicalSectorFactorization.zeroRightTensorOnInactive_neighboringOperator}
    \leanok
    \uses{def:mpdo_physical_sector_factorization, def:mpdo_minimal_neighboring_operator}
    Let $S$ be a set of active sectors in a physical-sector factorization,
    and suppose that $(l_k)_\beta=0$ for every $k\notin S$ and every $\beta$.
    Retain every physical summand and define
    \begin{align}
        \widehat l_k & =l_k,
        \notag               \\
        \widehat r_k & =
        \begin{cases}
            r_k, & k\in S,    \\
            0,   & k\notin S.
        \end{cases}
        \notag
    \end{align}
    Then $\widehat l_k\otimes\widehat r_k=l_k\otimes r_k$ for every $k$,
    so the physical-slice factorization is unchanged. The neighboring
    operators become
    \begin{align}
        \widehat\eta_{k,h}
         & =
        \begin{cases}
            \eta_{k,h}, & k\in S,    \\
            0,          & k\notin S.
        \end{cases}
        \notag
    \end{align}
    This is the reparameterization freedom in the factorization and
    neighboring contraction of
    \cite[Appendix~C.2, lines~1434--1445]{Cirac2016MPDO_arXiv}.
\end{definition}

\begin{definition}[Inverse-map physical-sector factorization]
    \label{def:mpdo_inverse_map_physical_sector_factorization}
    \lean{MPOTensor.inverseMapPhysicalSectorFactorization}
    \leanok
    \uses{thm:mpdo_sector_factorization, def:mpdo_physical_sector_factorization, lem:mpdo_trace_one_nonempty}
    Let $\mc K$ be injective, let $\rho$ be its three-site closure against
    a virtual tail $R$, and choose a Hayashi decomposition of $\rho$. Fix
    outer indices such that $R_{\beta_3,\alpha_1}\ne0$. The inverse-map
    tensors $l_k,r_k$, together with the Hayashi physical decomposition and
    unitary, determine a physical-sector factorization of $\mc K$.

    The left and right factor spaces are nonzero because the corresponding
    Hayashi density matrices have trace one; their nonzero dimensions are
    conclusions rather than hypotheses. No positivity of the neighboring
    operators is asserted. This gives the factorization in
    \cite[Appendix~C.2, equations \texttt{AppUkU=rl} and \texttt{formK},
        lines~1383--1387 and~1434--1439]{Cirac2016MPDO_arXiv}.
\end{definition}

\begin{theorem}[Identification of the neighboring operators]
    \label{thm:mpdo_inverse_map_factorization_neighboring_operator}
    \lean{MPOTensor.inverseMapPhysicalSectorFactorization_neighboringOperator_eq_sectorEta}
    \leanok
    \uses{def:mpdo_inverse_map_physical_sector_factorization, def:mpdo_sector_eta, def:mpdo_minimal_neighboring_operator}
    For all sectors $k,h$, the neighboring contraction of the resulting
    physical-sector factorization is the inverse-map operator
    \begin{align}
        \eta_{k,h}
         & =\sum_\gamma (r_k)_\gamma\otimes(l_h)_\gamma.
        \notag
    \end{align}
    This is equation \texttt{etarl} in
    \cite[Appendix~C.2, lines~1441--1445]{Cirac2016MPDO_arXiv}.
\end{theorem}

\begin{definition}[Zero-weight inverse-map reparameterization]
    \label{def:mpdo_zero_weight_reparameterized_inverse_map_factorization}
    \lean{MPOTensor.zeroWeightReparameterizedInverseMapPhysicalSectorFactorization}
    \leanok
    \uses{def:mpdo_inverse_map_physical_sector_factorization, def:mpdo_physical_sector_inactive_right_zero, lem:mpdo_sector_tensor_left_zero_of_weight_zero}
    In the inverse-map factorization, retain every physical sector and its
    factor spaces, and set
    \begin{align}
        \widehat l_k & =l_k,
        \notag               \\
        \widehat r_k & =
        \begin{cases}
            r_k, & p_k\ne0, \\
            0,   & p_k=0.
        \end{cases}
        \notag
    \end{align}
    The physical-slice factorization is unchanged. This treats the
    zero-weight freedom left by the Hayashi decomposition used in
    \cite[Appendix~C.2, lines~1415--1445]{Cirac2016MPDO_arXiv}.
\end{definition}

\begin{theorem}[Neighboring operators after zero-weight reparameterization]
    \label{thm:mpdo_zero_weight_reparameterized_neighboring_operator}
    \lean{MPOTensor.zeroWeightReparameterizedInverseMapPhysicalSectorFactorization_neighboringOperator}
    \lean{MPOTensor.probability_ne_zero_of_reparameterized_neighboringOperator_ne_zero}
    \leanok
    \uses{def:mpdo_zero_weight_reparameterized_inverse_map_factorization, thm:mpdo_inverse_map_factorization_neighboring_operator}
    The neighboring operators of the reparameterized factorization are
    \begin{align}
        \widehat\eta_{k,h}
         & =
        \begin{cases}
            \eta_{k,h}, & p_k\ne0, \\
            0,          & p_k=0.
        \end{cases}
        \notag
    \end{align}
    Thus a zero-weight sector has no outgoing support edge. Since
    $p_h=0$ also gives $l_h=0$, it has no incoming support edge either.
\end{theorem}

\begin{proof}
    \leanok
    The neighboring contraction uses $r_k$ at its source and $l_h$ at its
    target. The first assertion follows from the definition of
    $\widehat r_k$, and the last assertion follows from
    Lemma~\ref{lem:mpdo_sector_tensor_left_zero_of_weight_zero}.
\end{proof}

\begin{corollary}[SAL gives a raw physical-sector factorization]
    \label{cor:mpdo_sal_raw_physical_sector_factorization}
    \lean{MPOTensor.nonempty_physicalSectorFactorization_of_isSAL}
    \leanok
    \uses{thm:mpdo_four_site_sal_eta_structure, thm:mpdo_normalized_four_site_tail_entry_ne_zero, def:mpdo_inverse_map_physical_sector_factorization}
    Every injective MPO tensor satisfying SAL admits a physical-sector
    factorization. This is only the raw algebraic factorization from Lemma C.4
    and equations \texttt{formK} and \texttt{etarl}; coherent positivity of
    the neighboring operators, recurrence, and primitivity are not part of
    this conclusion.
\end{corollary}

\begin{definition}[Finite-chain coordinates of a physical-sector factorization]
    \label{def:mpdo_physical_sector_chain_coordinates}
    \lean{MPOTensor.PhysicalSectorFactorization.SectorChainFiber,
        MPOTensor.PhysicalSectorFactorization.SectorChainIndex,
        MPOTensor.PhysicalSectorFactorization.sectorChainEquiv,
        MPOTensor.PhysicalSectorFactorization.cyclicNeighboringProduct}
    \leanok
    \uses{def:mpdo_physical_sector_factorization, def:mpdo_minimal_neighboring_operator}
    Let $N\geq1$. For a cyclic sector configuration
    $k=(k_0,\ldots,k_{N-1})$, write
    \begin{align}
        I_k
         & =\prod_{n=0}^{N-1}(B_{k_n}^L\times B_{k_n}^R).
        \notag
    \end{align}
    Applying the one-site sector decomposition at every site gives a
    bijection
    \begin{align}
        \Phi_N:\Fin{d}^{\,N}
         & \longrightarrow
        \coprod_{k_0,\ldots,k_{N-1}}I_k.
        \notag
    \end{align}
    On $I_k$, define the cyclic neighboring product by
    \begin{align}
        \Omega_k(x,y)
         & =\prod_{n=0}^{N-1}
        \eta_{k_n,k_{n+1}}
        ((x_n^R,x_{n+1}^L),(y_n^R,y_{n+1}^L)),
        \notag
    \end{align}
    where site labels are read modulo $N$.
\end{definition}

\begin{theorem}[Finite-chain decomposition of a physical-sector factorization]
    \label{thm:mpdo_physical_sector_chain_decomposition}
    \lean{MPOTensor.PhysicalSectorFactorization.mpo_submatrix_sector_eq_cyclicNeighboringProduct,
        MPOTensor.PhysicalSectorFactorization.mpo_reindex_sectorChainEquiv_eq_blockDiagonal}
    \leanok
    \uses{def:mpdo_physical_sector_chain_coordinates}
    Suppose that the physical-sector factorization is given. For every
    $N\geq1$, the physically transformed MPO is block diagonal in the cyclic
    sector configurations, and
    \begin{align}
        \reindex_{\Phi_N}(M_N(\widetilde{\mc K}))
         & =\bigoplus_{k_0,\ldots,k_{N-1}}\Omega_k.
        \notag
    \end{align}
    This is the all-length algebraic decomposition in
    \cite[Appendix~C.2, lines~1435--1450]{Cirac2016MPDO_arXiv}, conditional
    here on the displayed physical-sector factorization. It does not assert
    positivity of the individual neighboring operators.
\end{theorem}

\begin{proof}
    \leanok
    Fix a sector configuration $k$ and entries $x,y\in I_k$. Expanding the
    closed horizontal contraction gives
    \begin{align}
        \sum_g\prod_n
        (l_{k_n})_{g_n}(x_n^L,y_n^L)
        (r_{k_n})_{g_{n+1}}(x_n^R,y_n^R).
        \notag
    \end{align}
    Cyclically shifting the summation variables and collecting the factors
    with the same horizontal index turns this expression into
    $\Omega_k(x,y)$. If the row and column sector configurations differ,
    one transformed physical slice lies between distinct direct summands and
    vanishes. Thus all off-diagonal sector blocks vanish.
\end{proof}

\begin{corollary}[SAL gives positive cyclic sector products]
    \label{cor:mpdo_sal_positive_cyclic_sector_products}
    \lean{MPOTensor.exists_physicalSectorFactorization_with_cyclicNeighboringProduct_posSemidef_of_isSAL}
    \leanok
    \uses{cor:mpdo_sal_raw_physical_sector_factorization, thm:mpdo_conjugated_chain_congruence, thm:mpdo_physical_sector_chain_decomposition}
    Every injective MPO tensor satisfying SAL admits a physical-sector
    factorization such that, for every $N\geq1$ and every cyclic sector
    configuration $k=(k_0,\ldots,k_{N-1})$ with $k_N=k_0$,
    \begin{align}
        \bigotimes_{n=0}^{N-1}\eta_{k_n,k_{n+1}}
         & \geq0.
        \notag
    \end{align}
    This is the projected-chain inequality in
    \cite[Appendix~C.2, Lemma~C.4, lines~1446--1450]{Cirac2016MPDO_arXiv}.
    It does not assert positivity of the individual neighboring operators.
\end{corollary}

\begin{proof}
    \leanok
    \uses{cor:mpdo_sal_raw_physical_sector_factorization, thm:mpdo_conjugated_chain_congruence, thm:mpdo_physical_sector_chain_decomposition}
    Choose the raw physical-sector factorization supplied by SAL. The physical
    basis change sends the positive $N$-site MPO to a positive congruence.
    Each cyclic neighboring product is a diagonal sector compression of this
    congruence and is therefore positive semidefinite.
\end{proof}

\begin{definition}[Boundary contraction against a virtual matrix]
    \label{def:mpdo_sector_boundary_operator}
    \lean{MPOTensor.PhysicalSectorFactorization.boundaryOperator,
        MPOTensor.PhysicalSectorFactorization.boundaryOperator_apply,
        MPOTensor.PhysicalSectorFactorization.boundaryOperator_zero,
        MPOTensor.PhysicalSectorFactorization.boundaryOperator_add,
        MPOTensor.PhysicalSectorFactorization.boundaryOperator_smul}
    \leanok
    \uses{def:mpdo_physical_sector_factorization}
    For an arbitrary virtual matrix $X$, define the outer boundary operator
    on $B_k^L\otimes B_h^R$ by
    \begin{align}
        B_{k,h}(X)
         & =\sum_{a,b}X_{b,a}\,(l_k)_a\otimes(r_h)_b.
        \notag
    \end{align}
    The map $X\mapsto B_{k,h}(X)$ is complex-linear. No positivity of $X$ or
    $B_{k,h}(X)$ is assumed.
\end{definition}

\begin{definition}[Regrouped two-site sector closure]
    \label{def:mpdo_two_site_sector_closure}
    \lean{MPOTensor.PhysicalSectorFactorization.twoSiteRegroupEquiv,
        MPOTensor.PhysicalSectorFactorization.twoSiteRegroupEquiv_symm_apply,
        MPOTensor.PhysicalSectorFactorization.twoSiteSectorEmbedding,
        MPOTensor.PhysicalSectorFactorization.twoSiteSectorClosure}
    \leanok
    \uses{def:phys_close, def:mpdo_physical_sector_factorization}
    First pass to the physical basis selected by $U$, so that each slice is
    written in the sector coordinates of
    Definition~\ref{def:mpdo_physical_sector_factorization}. Form the
    two-site closure of this sector-coordinate tensor, restrict it to the
    physical sectors $(k,h)$, and reindex its physical factors by
    \begin{align}
        ((L_k,R_k),(L_h,R_h))
         & \longmapsto ((L_k,R_h),(R_k,L_h)).
        \notag
    \end{align}
    Denote the resulting matrix by $\mc K_{2;k,h}(X)$.
\end{definition}

\begin{theorem}[Two-site closure factorization]
    \label{thm:mpdo_two_site_sector_closure_factorization}
    \lean{MPOTensor.PhysicalSectorFactorization.twoSiteSectorClosure_eq}
    \leanok
    \uses{def:mpdo_two_site_sector_closure, def:mpdo_minimal_neighboring_operator, def:mpdo_sector_boundary_operator}
    For every virtual matrix $X$ and every sector pair $(k,h)$,
    \begin{align}
        \mc K_{2;k,h}(X)
         & =B_{k,h}(X)\otimes\eta_{k,h}.
        \notag
    \end{align}
\end{theorem}
\begin{proof}
    \leanok
    \uses{def:mpdo_two_site_sector_closure, def:mpdo_minimal_neighboring_operator, def:mpdo_sector_boundary_operator, def:phys_close, def:mpdo_physical_sector_factorization}
    At the regrouped matrix entry indexed by
    $((\lambda_k,\rho_h),(\rho_k,\lambda_h))$ and
    $((\lambda'_k,\rho'_h),(\rho'_k,\lambda'_h))$, the virtual trace expands
    as the scalar
    \begin{align}
        \sum_{a,b,c}X_{b,a}
        [(l_k)_a]_{\lambda_k,\lambda'_k}
        [(r_k)_c]_{\rho_k,\rho'_k}
        [(l_h)_c]_{\lambda_h,\lambda'_h}
        [(r_h)_b]_{\rho_h,\rho'_h}.
        \notag
    \end{align}
    Commuting scalar factors and collecting the sums over $(a,b)$ and $c$
    gives respectively $B_{k,h}(X)$ and $\eta_{k,h}$.
\end{proof}

\begin{definition}[Regrouped three-site sector closure]
    \label{def:mpdo_three_site_sector_closure}
    \lean{MPOTensor.PhysicalSectorFactorization.threeSiteRegroupEquiv,
        MPOTensor.PhysicalSectorFactorization.threeSiteRegroupEquiv_symm_apply,
        MPOTensor.PhysicalSectorFactorization.threeSiteSectorEmbedding,
        MPOTensor.PhysicalSectorFactorization.threeSiteSectorClosure}
    \leanok
    \uses{def:phys_close, def:mpdo_physical_sector_factorization}
    First pass to the physical basis selected by $U$ and form the three-site
    closure of the resulting sector-coordinate tensor. Restrict this closure
    to the sectors $(k,l,h)$ and reindex its physical factors by
    \begin{align}
        ((L_k,R_k),((L_l,R_l),(L_h,R_h)))
         & \longmapsto
        ((L_k,R_h),((R_k,L_l),(R_l,L_h))).
        \notag
    \end{align}
    Denote the resulting matrix by $\mc K_{3;k,l,h}(X)$.
\end{definition}

\begin{theorem}[Three-site closure factorization]
    \label{thm:mpdo_three_site_sector_closure_factorization}
    \lean{MPOTensor.PhysicalSectorFactorization.threeSiteSectorClosure_eq}
    \leanok
    \uses{def:mpdo_three_site_sector_closure, def:mpdo_minimal_neighboring_operator, def:mpdo_sector_boundary_operator}
    For every virtual matrix $X$ and fixed sectors $(k,l,h)$,
    \begin{align}
        \mc K_{3;k,l,h}(X)
         & =B_{k,h}(X)\otimes
        (\eta_{k,l}\otimes\eta_{l,h}).
        \notag
    \end{align}
    \begin{center}
        \tenkzkernel
        % K_{3;k,l,h}(X) = B_{k,h}(X) (x) eta_{k,l} (x) eta_{l,h}: the
        % three-site closure of the sector-restricted MPO chain, after
        % regrouping its six bond-space legs (L_k,R_k,L_l,R_l,L_h,R_h),
        % separates into the boundary factor and the two neighboring
        % factors. Ink-to-index correspondence: l_k, r_k, l_l, r_l,
        % l_h, r_h are the sector-coordinate tensors of
        % Definition~\ref{def:mpdo_physical_sector_factorization}; X is
        % the virtual matrix of the physical closure of
        % Definition~\ref{def:phys_close}. Every open leg (up or down)
        % is one row or column index of a bond space, named at its
        % tip. Both sides carry six open leg-pairs, grouped left to
        % right as $(L_k,R_h)\mid(R_k,L_l)\mid(R_l,L_h)$, matching the
        % regrouping map of
        % Definition~\ref{def:mpdo_three_site_sector_closure}
        % (virtual-west=0 | virtual-east=0 | physical-up=6 |
        % physical-down=6 on both sides).
        \begin{tenkz}[rows={wire}, cols=6, west=none, east=none]
            \tn[skin=pill, wide=6,
                ports={90@1:physical:$L_k$, 90@2:physical:$R_h$,
                        90@3:physical:$R_k$, 90@4:physical:$L_l$,
                        90@5:physical:$R_l$, 90@6:physical:$L_h$,
                        270@1:physical:$L_k$, 270@2:physical:$R_h$,
                        270@3:physical:$R_k$, 270@4:physical:$L_l$,
                        270@5:physical:$R_l$, 270@6:physical:$L_h$}]
            {{\cal K}_{3;k,l,h}(X)}
        \end{tenkz}
        $\;=\;$
        % B_{k,h}(X): l_k and r_h flank the virtual matrix X on the
        % wire; their open legs are L_k and R_h.
        \begin{tenkz}[rows={wire}, cols=3, west=none, east=none]
            \tn[label pos=225, ports={90:physical:$L_k$, 270:physical:$L_k$}]{l_k} &
            \tn[skin=ring]{X} &
            \tn[label pos=225, ports={90:physical:$R_h$, 270:physical:$R_h$}]{r_h}
        \end{tenkz}
        $\;\otimes\;$
        % eta_{k,l}: r_k and l_l share the contracted virtual bond
        % between sites k and l; their open legs are R_k and L_l.
        \begin{tenkz}[rows={wire}, cols=2, west=none, east=none]
            \tn[label pos=225, ports={90:physical:$R_k$, 270:physical:$R_k$}]{r_k} &
            \tn[label pos=225, ports={90:physical:$L_l$, 270:physical:$L_l$}]{l_l}
        \end{tenkz}
        $\;\otimes\;$
        % eta_{l,h}: r_l and l_h share the contracted virtual bond
        % between sites l and h; their open legs are R_l and L_h.
        \begin{tenkz}[rows={wire}, cols=2, west=none, east=none]
            \tn[label pos=225, ports={90:physical:$R_l$, 270:physical:$R_l$}]{r_l} &
            \tn[label pos=225, ports={90:physical:$L_h$, 270:physical:$L_h$}]{l_h}
        \end{tenkz}
    \end{center}
\end{theorem}
\begin{proof}
    \leanok
    \uses{def:mpdo_three_site_sector_closure, def:mpdo_minimal_neighboring_operator, def:mpdo_sector_boundary_operator, def:phys_close, def:mpdo_physical_sector_factorization}
    At the regrouped matrix entry indexed by
    $((\lambda_k,\rho_h),((\rho_k,\lambda_l),(\rho_l,\lambda_h)))$ and its
    primed counterpart, the left-hand side is the scalar
    \begin{align}
        \sum_{a,b,c,e}X_{b,a}
        [(l_k)_a]_{\lambda_k,\lambda'_k}
        [(r_k)_c]_{\rho_k,\rho'_k}
        [(l_l)_c]_{\lambda_l,\lambda'_l}
        [(r_l)_e]_{\rho_l,\rho'_l}
        [(l_h)_e]_{\lambda_h,\lambda'_h}
        [(r_h)_b]_{\rho_h,\rho'_h}.
        \notag
    \end{align}
    Reordering the four finite sums separates the $(a,b)$ boundary
    contraction from the neighboring contractions indexed by $c$ and $e$,
    which gives the three factors on the right-hand side.
\end{proof}

\begin{definition}[Three-site neighboring operator]
    \label{def:mpdo_physical_sector_three_site_neighboring_operator}
    \lean{MPOTensor.PhysicalSectorFactorization.threeSiteNeighboringOperator}
    \leanok
    \uses{def:mpdo_minimal_neighboring_operator}
    For every outer-sector pair $(k,h)$, define
    \begin{align}
        \Omega_{k,h}
         & =\bigoplus_l
        (\eta_{k,l}\otimes\eta_{l,h}).
        \notag
    \end{align}
    This is the direct sum of the two neighboring contractions appearing in
    the three-site closure factorization
    \cite[Appendix~C.2, lines~1510--1516]{Cirac2016MPDO_arXiv}.
\end{definition}

\begin{figure}[ht]
    \centering
    \begin{minipage}{0.41\textwidth}
        \centering
        \begin{align}
            \mathfrak R_2(\mc K_2(X))_{k,h}
             & =a_kb_hB_{k,h}(X).
            \notag
        \end{align}
        \begin{center}
            \tenkzkernel
            % R_2(K_2(X))_{k,h} = a_kb_h B_{k,h}(X): the two-site
            % closure, restricted to sector pair (k,h) and traced over
            % its inner bond space R_k (x) L_h, factors into the scalar
            % a_kb_h = tr(eta_{k,h}) times the boundary operator
            % B_{k,h}(X) = sum_{a,b} X_{b,a} (l_k)_a (x) (r_h)_b.
            % Ink-to-index correspondence: the left pill is the
            % already-traced closure, its up/down legs the row/column
            % indices of L_k, R_h; on the right, l_k and r_h are the
            % sector-coordinate tensors of
            % Definition~\ref{def:mpdo_sector_boundary_operator}, X the
            % virtual matrix riding the bond between their own
            % horizontal indices a, b, and its two ends are sealed --
            % no further horizontal index survives. Both sides carry
            % the same boundary signature (virtual-west=0 |
            % virtual-east=0 | physical-up=2 | physical-down=2).
            \begin{tenkz}[rows={wire}, cols=2, west=none, east=none]
                \tn[skin=pill, wide=2,
                    ports={90@1:physical:$L_k$, 90@2:physical:$R_h$,
                            270@1:physical:$L_k$, 270@2:physical:$R_h$}]
                {\mathfrak R_2({\cal K}_2(X))_{k,h}}
            \end{tenkz}
            $\;=\;a_kb_h\;$
            \begin{tenkz}[rows={wire}, cols=3, west=none, east=none]
                \tn[label pos=225, ports={90:physical:$L_k$, 270:physical:$L_k$}]{l_k} &
                \tn[skin=ring]{X} &
                \tn[label pos=225, ports={90:physical:$R_h$, 270:physical:$R_h$}]{r_h}
            \end{tenkz}
        \end{center}
    \end{minipage}\hfill
    \begin{minipage}{0.56\textwidth}
        \centering
        \begin{align}
            \mathfrak R_3(\mc K_3(X))_{k,l,h}
             & =B_{k,h}(X)\otimes\eta_{k,l}\otimes\eta_{l,h}.
            \notag
        \end{align}
    \end{minipage}
    \caption{First each slice is conjugated by $U$ and expressed in the
        corresponding sector coordinates. After restriction to fixed sectors and
        the displayed regroupings, the arbitrary-$X$ two- and three-site closures
        separate into the boundary factor and respectively one or two neighboring
        contractions. The displayed identities do not assert positivity,
        preparation, recovery, or the full conclusion of Proposition~C.7. They
        follow the contractions in
        \cite[lines~638--655 and Appendix~C.2, lines~1435--1448]
        {Cirac2016MPDO_arXiv}.}
    \label{fig:mpdo_sector_closure_factorizations}
\end{figure}

\begin{definition}[Neighboring-operator trace factorization]
    \label{def:mpdo_neighboring_trace_factorization}
    \lean{MPOTensor.PhysicalSectorFactorization.NeighboringTraceFactorization}
    \lean{MPOTensor.PhysicalSectorFactorization.NeighboringTraceFactorization.weight_nonneg,
        MPOTensor.PhysicalSectorFactorization.NeighboringTraceFactorization.neighboringOperator_eq_zero_of_mul_eq_zero,
        MPOTensor.PhysicalSectorFactorization.NeighboringTraceFactorization.neighborIndex_nonempty_of_mul_ne_zero}
    \leanok
    \uses{def:mpdo_minimal_neighboring_operator}
    Suppose that every neighboring operator is positive semidefinite and that
    there are real numbers $a_k,b_h$ such that
    \begin{align}
        \tr(\eta_{k,h}) & =a_kb_h,
        \notag                     \\
        \sum_k a_kb_k   & =1.
        \notag
    \end{align}
    These are precisely the remaining hypotheses used in the construction of
    Proposition~C.7 after the physical-sector factorization has been fixed
    \cite[Appendix~C.2, lines~1389--1403]{Cirac2016MPDO_arXiv}.
\end{definition}

% This coordinate-invariance theorem is project-derived.  It transports the
% data in CPSV16 Theorem 4.9(iv), but is not a theorem stated in that source.
\begin{theorem}[Unitary-between physical-coordinate invariance of neighboring trace factorizations]
    \label{thm:mpdo_neighboring_trace_factorization_physical_basis_invariant}
    \lean{MPOTensor.PhysicalSectorFactorization.ofChangePhysicalBasis,
        MPOTensor.PhysicalSectorFactorization.toChangePhysicalBasis,
        MPOTensor.PhysicalSectorFactorization.ofChangePhysicalBasis_neighboringOperator,
        MPOTensor.PhysicalSectorFactorization.toChangePhysicalBasis_neighboringOperator,
        MPOTensor.PhysicalSectorFactorization.NeighboringTraceFactorization.ofChangePhysicalBasis,
        MPOTensor.PhysicalSectorFactorization.NeighboringTraceFactorization.toChangePhysicalBasis,
        MPOTensor.PhysicalSectorFactorization.exists_neighboringTraceFactorization_changePhysicalBasis_iff,
        MPOTensor.PhysicalSectorFactorization.exists_neighboringTraceFactorization_changePhysicalBasis_iff_of_isUnitaryBetween}
    \leanok
    \uses{def:mpdo_physical_sector_factorization, def:mpdo_neighboring_trace_factorization}
    Let $V:\C^d\to\C^e$ satisfy
    $V^\dagger V=\Id_d$ and $VV^\dagger=\Id_e$, and define
    ${\cal K}^V_{\beta,\alpha}
    =V{\cal K}_{\beta,\alpha}V^\dagger$.  Then ${\cal K}$ admits a
    physical-sector factorization with
    \begin{align}
        \eta_{k,h} & \geq 0,
        \notag \\
        \tr(\eta_{k,h}) & =a_kb_h,
        \notag \\
        \sum_k a_kb_k & =1.
        \notag
    \end{align}
    if and only if ${\cal K}^V$ admits one.  The same sector spaces, tensors
    $l_k,r_k$, neighboring operators $\eta_{k,h}$, and numbers $a_k,b_h$ may
    be used on both sides.  This invariance is project-derived; the data being
    transported are those of
    \cite[Theorem~4.9(iv) and Appendix~C.2, lines~1381--1450]
    {Cirac2016MPDO_arXiv}, but the equivalence itself is not stated there.
\end{theorem}

\begin{proof}
    \leanok
    The two one-sided unitary identities give $d\leq e$ and $e\leq d$,
    respectively, so $d=e$.  If the physical isometry for
    ${\cal K}$ is $U$, use $UV^\dagger$ for ${\cal K}^V$; then
    $(UV^\dagger){\cal K}^V_{\beta,\alpha}(UV^\dagger)^\dagger
    =U{\cal K}_{\beta,\alpha}U^\dagger$.  Conversely, if the isometry for
    ${\cal K}^V$ is $U'$, use $U'V$ for ${\cal K}$.  In both directions the
    tensors $l_k,r_k$ are unchanged, so
    $\eta_{k,h}=\sum_\gamma(r_k)_\gamma\otimes(l_h)_\gamma$ is unchanged.
    The positivity, trace, and normalization identities therefore remain the
    same identities.
\end{proof}

% These presentation-independence statements are project-derived.  CPSV16
% supplies the transported data and the source settings for gauge and support
% restriction, but does not state either equivalence below.
\begin{theorem}[Gauge and support-coordinate independence of neighboring trace factorizations]
    \label{thm:mpdo_neighboring_trace_factorization_presentation_invariant}
    \lean{MPOTensor.PhysicalSectorFactorization.NeighboringTraceFactorization.ofGaugeEquiv,
        MPOTensor.PhysicalSectorFactorization.exists_neighboringTraceFactorization_iff_of_gaugeEquiv,
        MPOTensor.PhysicalSupportRestrictionData.supportCoordinateChange,
        MPOTensor.PhysicalSupportRestrictionData.supportCoordinateChange_isUnitaryBetween,
        MPOTensor.PhysicalSupportRestrictionData.changePhysicalBasis_supportCoordinateChange,
        MPOTensor.PhysicalSupportRestrictionData.exists_neighboringTraceFactorization_restriction_iff}
    \leanok
    \uses{def:gauge_equiv, def:mpdo_neighboring_trace_factorization,
        def:mpdo_physical_support_restriction}
    Suppose that two MPO tensors of the same physical and virtual dimensions
    differ by an invertible virtual gauge.  One admits a physical-sector
    factorization with neighboring data $\eta_{k,h}$, $a_k$, and $b_h$ if and
    only if the other does.  The same neighboring operators and the same
    numbers $a_k,b_h$ may be used on both sides.

    Let $V_F:\C^{d_F}\to\C^d$ and
    $V_G:\C^{d_G}\to\C^d$ be two isometric inclusions with the same range
    projection $P$:
    \begin{align}
        V_F^\dagger V_F &=\Id_{d_F},
        & V_G^\dagger V_G &=\Id_{d_G},
        \notag \\
        V_FV_F^\dagger &=P,
        & V_GV_G^\dagger &=P.
        \label{eq:mpdo_support_coordinate_ranges}
    \end{align}
    Set $Q=V_F^\dagger V_G$.  Then
    \begin{align}
        Q^\dagger Q &=\Id_{d_G},
        & QQ^\dagger &=\Id_{d_F},
        \notag \\
        Q\bigl(V_G^\dagger {\cal K}V_G\bigr)Q^\dagger
          &=V_F^\dagger {\cal K}V_F.
        \label{eq:mpdo_support_coordinate_tensor}
    \end{align}
    Consequently, the restrictions of ${\cal K}$ along $V_F$ and $V_G$
    admit neighboring trace factorizations simultaneously.

    Both equivalences are project-derived and are not stated in
    \cite{Cirac2016MPDO_arXiv}.  The gauge transformation is the one discussed
    in Section~2.3, lines~187--195; the support-restriction setting is that of
    Appendix~C.2, lines~1680--1691 and 1733--1770.
\end{theorem}

\begin{proof}
    \leanok
    \uses{thm:mpdo_physical_sector_factorization_gauge,
        thm:mpdo_neighboring_trace_factorization_physical_basis_invariant}
    For the virtual gauge $L^i=XK^iX^{-1}$, write
    $\eta^K_{k,h}(M)$ for the neighboring contraction with $M$ inserted
    between the transported virtual factors.  Their composition is
    $X^{-1}X$, and hence
    \begin{align}
        \eta^L_{k,h}
        &=\eta^K_{k,h}(X^{-1}X)
         =\eta^K_{k,h}(\Id)
         =\eta^K_{k,h}.
        \notag
    \end{align}
    The families $a_k$ and $b_h$ are retained literally, together with their
    positivity, trace, and normalization identities.  The reverse implication
    uses the inverse gauge.

    For the support coordinates,
    (\ref{eq:mpdo_support_coordinate_ranges}) gives
    \begin{align}
        Q^\dagger Q
          &=V_G^\dagger(V_FV_F^\dagger)V_G
           =V_G^\dagger PV_G
           =\Id_{d_G},
        \notag \\
        QQ^\dagger
          &=V_F^\dagger(V_GV_G^\dagger)V_F
           =V_F^\dagger PV_F
           =\Id_{d_F},
        \notag \\
        QV_G^\dagger
          &=V_F^\dagger(V_GV_G^\dagger)
           =V_F^\dagger P
           =V_F^\dagger.
        \notag
    \end{align}
    The last identity gives
    (\ref{eq:mpdo_support_coordinate_tensor}) by composition of physical
    coordinate changes.  The unitary-between invariance above, with the
    resulting equivalence reversed, proves the claim.
\end{proof}

\begin{theorem}[Trace of the three-site neighboring operator]
    \label{thm:mpdo_physical_sector_three_site_neighboring_operator_trace}
    \lean{MPOTensor.PhysicalSectorFactorization.trace_threeSiteNeighboringOperator,
        MPOTensor.PhysicalSectorFactorization.NeighboringTraceFactorization.trace_threeSiteNeighboringOperator}
    \leanok
    \uses{def:mpdo_physical_sector_three_site_neighboring_operator, def:mpdo_neighboring_trace_factorization}
    Assume the neighboring-operator trace factorization of
    Definition~\ref{def:mpdo_neighboring_trace_factorization}. Then, for every
    outer-sector pair $(k,h)$,
    \begin{align}
        \tr(\Omega_{k,h})
         & =\sum_l\tr(\eta_{k,l})\tr(\eta_{l,h})
        =a_kb_h.
        \notag
    \end{align}
    This is the normalization used in the coarse-graining channel of
    \cite[Appendix~C.2]{Cirac2016MPDO_arXiv}, lines~1547--1555.
\end{theorem}

\begin{proof}
    \leanok
    \uses{def:mpdo_neighboring_trace_factorization}
    The trace of a direct sum is the sum of the traces of its blocks, and the
    trace of a tensor product is the product of the traces. Hence
    \begin{align}
        \tr(\Omega_{k,h})
         & =\sum_l(a_kb_l)(a_lb_h)
        =a_kb_h\sum_l a_lb_l
        =a_kb_h.
        \notag
    \end{align}
\end{proof}

\begin{lemma}[Restricting finite sums to nonzero weights]
    \label{lem:finset_sum_subtype_ne_zero}
    \lean{Finset.sum_eq_sum_subtype_ne_zero}
    \leanok
    Let $I$ be finite, let $p_i$ be weights in a type with zero, and let $f_i$
    take values in a commutative additive monoid. If $f_i=0$ whenever $p_i=0$,
    then
    \begin{align}
        \sum_{i\in I}f_i&=\sum_{\substack{i\in I\\p_i\ne0}}f_i.
        \notag
    \end{align}
\end{lemma}

\begin{proof}
    \leanok
    Split the full sum according to whether $p_i$ vanishes. Every term in the
    zero-weight part vanishes by hypothesis.
\end{proof}

\begin{lemma}[Cyclic trace identity for rectangular powers]
    \label{lem:matrix_trace_pow_mul_comm}
    \lean{Matrix.trace_pow_mul_comm}
    \leanok
    For rectangular matrices $L$ and $Q$ over a commutative semiring and every
    integer $N\geq1$,
    \begin{align}
        \tr\bigl((LQ)^N\bigr)&=\tr\bigl((QL)^N\bigr).
        \notag
    \end{align}
\end{lemma}

\begin{proof}
    \leanok
    Write $N=M+1$. Induction on $M$ gives
    $(LQ)^{M+1}=L(QL)^M Q$. Moving the first factor cyclically through the
    trace and reassociating the product gives $\tr((QL)^{M+1})$.
\end{proof}

\begin{theorem}[Unit closure traces from a neighboring trace factorization]
    \label{thm:mpdo_neighboring_trace_unit_closure}
    \lean{MPOTensor.PhysicalSectorFactorization.leftTraceMatrix,
      MPOTensor.PhysicalSectorFactorization.rightTraceMatrix,
      MPOTensor.PhysicalSectorFactorization.trace_neighboringOperator_eq_sum_trace_mul_trace,
      MPOTensor.PhysicalSectorFactorization.physTraceTransfer_eq_leftTraceMatrix_mul_rightTraceMatrix,
      MPOTensor.PhysicalSectorFactorization.rightTraceMatrix_mul_leftTraceMatrix_apply,
      MPOTensor.PhysicalSectorFactorization.NeighboringTraceFactorization.rightTraceMatrix_mul_leftTraceMatrix_eq_vecMulVec,
      MPOTensor.PhysicalSectorFactorization.NeighboringTraceFactorization.trace_physTraceTransfer_pow_eq_one,
      MPOTensor.PhysicalSectorFactorization.NeighboringTraceFactorization.trace_physTraceTransfer_eq_one,
      MPOTensor.PhysicalSectorFactorization.NeighboringTraceFactorization.trace_mpo_eq_one}
    \leanok
    \uses{def:mpdo_physical_sector_factorization,
      def:mpdo_neighboring_trace_factorization,
      def:mpo_physical_trace_transfer, def:mpo_family}
    Let $M$ be an MPO tensor admitting a physical-sector factorization whose
    neighboring operators satisfy the trace factorization of
    Definition~\ref{def:mpdo_neighboring_trace_factorization}. Closing the
    physical legs of the factorized slices writes the physical-trace transfer
    as a rectangular product
    \begin{align}
        \mathcal T_M&=LQ,
        &L_{\beta,k}&=\tr\bigl((l_k)_\beta\bigr),
        &Q_{k,\alpha}&=\tr\bigl((r_k)_\alpha\bigr),
        \notag
    \end{align}
    whose opposite product is the rank-one coefficient matrix
    \begin{align}
        (QL)_{k,h}&=\tr(\eta_{k,h})=a_kb_h.
        \notag
    \end{align}
    Consequently every positive power of the physical-trace transfer and
    every nonempty closed chain has unit trace:
    \begin{align}
        \tr\bigl(\mathcal T_M^{\,N}\bigr)&=1,
        &\tr\rho^{(N)}(M)&=1,
        &N&\geq1.
        \notag
    \end{align}
    In particular, a tensor whose closed chains do not all have unit trace
    admits no physical-sector factorization with positive neighboring
    operators and normalized rank-one trace coefficients. The normalization
    $\tr(T)=1$ invoked in
    \cite[Appendix~C.2, line~1498]{Cirac2016MPDO_arXiv} is therefore a
    constraint on the tensor rather than a convention, and the ambient
    normalization at
    \cite[Appendix~C.2, lines~1755--1759]{Cirac2016MPDO_arXiv} constrains
    only the weighted sum over the representatives of a basis of normal
    tensors, not each representative separately.
\end{theorem}

\begin{proof}
    \leanok
    \uses{def:mpdo_neighboring_trace_factorization,
      def:mpo_physical_trace_transfer, thm:mpdo_closed_sector_physical_trace_transfer,
      lem:matrix_trace_pow_mul_comm, lem:mpdo_trace_loop_chain}
    Closing the physical legs of each factorized slice and using unitary
    invariance of the trace expresses each entry of the physical-trace
    transfer as $\sum_k\tr((l_k)_\beta)\tr((r_k)_\alpha)$, which is the
    rectangular product $LQ$. The entries of the opposite product are the
    contractions
    \begin{align}
        (QL)_{k,h}
        &=\sum_\alpha\tr\bigl((r_k)_\alpha\bigr)\tr\bigl((l_h)_\alpha\bigr)
         =\tr(\eta_{k,h})=a_kb_h.
        \notag
    \end{align}
    Cyclic invariance of the trace gives
    $\tr((LQ)^N)=\tr((QL)^N)$ for every $N\geq1$, and every positive power
    of the rank-one matrix $(a_kb_h)_{k,h}$ with $\sum_ka_kb_k=1$ equals the
    matrix itself, of trace one. The closed chain of length $N$ has trace
    $\tr(\mathcal T_M^{\,N})$.
\end{proof}

\begin{theorem}[Partial traces of the sector closures]
    \label{thm:mpdo_sector_closure_partial_traces}
    \lean{MPOTensor.PhysicalSectorFactorization.partialTraceRight_twoSiteSectorClosure}
    \lean{MPOTensor.PhysicalSectorFactorization.NeighboringTraceFactorization.partialTraceRight_twoSiteSectorClosure}
    \lean{MPOTensor.PhysicalSectorFactorization.partialTraceRight_threeSiteSectorClosure}
    \lean{MPOTensor.PhysicalSectorFactorization.NeighboringTraceFactorization.sum_threeSite_trace_coefficients}
    \leanok
    \uses{thm:mpdo_two_site_sector_closure_factorization, thm:mpdo_three_site_sector_closure_factorization, def:mpdo_neighboring_trace_factorization}
    For every virtual matrix $X$,
    \begin{align}
        \tr_{R_k\otimes L_h}\!\left(\mc K_{2;k,h}(X)\right)
         & =\tr(\eta_{k,h})B_{k,h}(X)
        =a_kb_h B_{k,h}(X).
        \notag
    \end{align}
    Similarly,
    \begin{align}
        \tr_{(R_k\otimes L_l)\otimes(R_l\otimes L_h)}
        \!\left(\mc K_{3;k,l,h}(X)\right)
         & =\tr(\eta_{k,l})\tr(\eta_{l,h})B_{k,h}(X),
        \notag
    \end{align}
    and summing the scalar coefficient over $l$ gives $a_kb_h$.
\end{theorem}

\begin{proof}
    \leanok
    \uses{thm:mpdo_two_site_sector_closure_factorization, thm:mpdo_three_site_sector_closure_factorization}
    The partial trace of $A\otimes B$ is $\tr(B)A$. The two-site identity
    follows immediately. For three sites, the coefficient is
    \begin{align}
        \sum_l (a_kb_l)(a_lb_h)
         & =a_k\left(\sum_l a_lb_l\right)b_h
        =a_kb_h.
        \notag
    \end{align}
\end{proof}
